\honest{Lost Purposes}{Eliezer Yudkowsky}{2007}

In kindergarten or first grade I was told that a Hebrew prayer works even if you do not know what the words mean. That was the beginning of my break with Judaism.

Students study ``12th-century knitting patterns'' because a high-paying job needs a credential, which needs a master's, which needs a bachelor's, which needs the class; they study hard meaning to forget it all after the exam. Maybe you saw it was madness and did it anyway. In the Bay Area, 80\% of elementary teachers spend less than an hour a week on science, and 16\% none. I understand the cause to be the No Child Left Behind Act. ``Virtually all classroom time is now spent on preparing for tests,'' and I seem to recall, without a source, that taking them filled 40\% of classroom time in one school. \nb{The survey's figures are as the post reports them. The ``virtually all'' has no source. A national survey the same year found elementary schools still spending about three hours a week on science, and 44\% of districts cutting other subjects to add reading and mathematics.}

Soviet factories met their quotas with tiny shoes, or by counting cut but unassembled leather as shoes, and their superiors, who also wanted to report success, did not look closely. It is ``now being suggested'' that most significant findings in medicine are untrue. But while $p<0.05$ is the bar for publication, why spend more on larger studies? Everyone knows the whole point of science is to publish papers, as the point of a university is to print parchment and of a school to pass the tests that fund it, while physics journals ``require a threshold of p<0.0001,'' as if they had some other purpose. \nb{The medical claim follows its source, a 2005 paper by Ioannidis. No such rule for physics journals was found. Particle physics requires a far smaller p-value, but only to claim a discovery.}

Nobody keeps opening the car door after learning the supermarket has no chocolate. In large organizations, though, you see what would be insanity in one mind: someone paid for every car door opened, who does not care whether anyone reaches the supermarket. To a Bayesian, subgoals are by-products of the probability function. There is no expected utility without utility. I trace No Child Left Behind through its links. Politicians must look busy to voters this year, not in fifteen years. Bureaucrats must show progress measurable this year. Textbook committees compare books by how many subjects they cover, and grades do not coordinate, so publishers cram in subjects, and teachers ``won't get through a fourth of the textbook.'' Teachers might complain, but do not decide. The consumers are the children, who cannot pay, vote or sit on committees; their parents can judge only surface images. Want it solved? Make the politicians go to school. I give no sources for any link. Schools ``seem'' to be ``collapsing into a state not much better than nothing.'' \nb{The US national assessment reported in 2007 that fourth-graders scored higher in reading and mathematics than in any earlier year.}

One mind can track expected utility through a dozen events, including a door whose value depends on the chocolate; organizations can reward only what is measurable and contractible today, which means intermediate events, and those are leaky generalizations. Bureaucrats are untrustworthy genies, for they do not share the wisher's values. \nb{The point that measured targets corrupt the process they measure, above all test scores in schools, is Donald Campbell's, from 1976. The post does not mention him.}

Musashi says that whenever you parry, strike or touch the enemy's sword, you must cut the enemy in the same movement, and ``You must thoroughly research this.'' A fighter taught by others, who did not generate the art, may not know why to parry now and spring then, or when the rule fails. Every rule must cut through to the truth, and every expected utility to utility. C. J. Cherryh: ``Your sword has no blade. It has only your intention.'' People who draft wishes for an imagined AI do not ask, as I ``reflexively'' ask, why they think a wish is good and whether the genie would judge likewise; they do not see the criterion behind their judgment. Nor do people notice selfish people giving altruistic arguments for selfishness, or the reverse. People track goals well inside their own heads, but dozens of organizations and years lie between a bored child and an incompetent graduate, and every link loses information and incentive. This bothers most people less than it bothers me, ``or why were all my classmates willing to say prayers without knowing what they meant?'' \nb{The classmates were small children.}

Can people learn to keep their eye on the ball? People do often want to do their jobs; can there be a sane corporation, or a sane civilization? That is what these posts on expected utility and utility have been aiming at. The worst threat to a complex civilization is its own complexity. My life ``has been driven by an exceptionally strong abhorrence to lost purposes,'' and I hope that can be taught.
